\section{A certified compiler for polynomial implicit equations}\label{sec:polynomial}

The gain theorem becomes a producer when its hypotheses can be computed from the input syntax. Coded rational polynomials provide a broad, nontrivial class for which this is possible. This section gives the missing construction rather than treating a Lipschitz estimate as an unspecified oracle.

\subsection{Finite difference bounds from coefficients}

Let \(d\ge1\), and let a multi-index \(\alpha=(\alpha_0,\ldots,\alpha_{d-1})\) be a finite list of naturals. Put \(x^\alpha=\prod_{j<d}x_j^{\alpha_j}\). For a rational box with positive radii \(R_j\), define
\[
 D_{\alpha,j}(R)=
 \begin{cases}
 \alpha_j R_j^{\alpha_j-1}\displaystyle\prod_{k\ne j}R_k^{\alpha_k},&\alpha_j>0,\\
 0,&\alpha_j=0.
 \end{cases}
\]
Every exponent in an executed branch is nonnegative.

\begin{lemma}[Monomial difference]\label{lem:monomial}
If \(\abs{x_j},\abs{y_j}\le R_j\), then
\[
 \abs{x^\alpha-y^\alpha}
 \le\sum_{j<d}D_{\alpha,j}(R)\abs{x_j-y_j}.
\]
\end{lemma}
\begin{proof}
For each positive integer \(m\), the finite identity
\(s^m-t^m=(s-t)\sum_{i<m}s^{m-1-i}t^i\)
gives \(\abs{s^m-t^m}\le mR^{m-1}\abs{s-t}\). For the product of the \(d\) powers, replace the factors one at a time. The difference between consecutive products is the difference in the replaced factor times unchanged factors, each bounded by the appropriate power of \(R_k\). Summing the resulting \(d\) inequalities proves the claim. Both identities are proved by induction on a finite product length. No derivative or mean-value theorem is used.
\end{proof}

Suppose \(T:\Q^d\to\Q^d\) is given by finite coefficient lists
\(T_i(x)=\sum_\alpha t_{i,\alpha}x^\alpha\). Define the nonnegative rational matrix
\begin{equation}\label{eq:polygain}
 W_{ij}=\sum_\alpha\abs{t_{i,\alpha}}D_{\alpha,j}(R).
\end{equation}
Then the preceding lemma gives the componentwise finite estimate
\begin{equation}\label{eq:polylip}
 \abs{T_i(x)-T_i(y)}\le\sum_{j<d}W_{ij}\abs{x_j-y_j}
\end{equation}
throughout the box. Computing \(W\) requires only the finite polynomial code and rational arithmetic. It may overestimate the true gain because it ignores cancellation. Failure of the certificate therefore means only that this certificate did not establish contraction.

\subsection{The finite certificate and its checker}

Take a polynomial system \(F:\Q^d\to\Q^d\), a rational center \(x_0\), matrices \(B,C\), a positive rational \(r\), and positive rational weights \(v_i\). Define
\begin{equation}\label{eq:polyT}
 T(h)=h-BF(x_0+h),\qquad R_i=r v_i.
\end{equation}
The finite certificate checker verifies dimensions and positivity, the identity \(CB=I\), and the inequalities
\begin{equation}\label{eq:polycert}
 2\abs{T_i(0)}\le r v_i,\qquad
 2\sum_{j<d}W_{ij}v_j\le v_i\quad(i<d),
\end{equation}
where \(W\) is computed from the expanded polynomial \(T\) by \eqref{eq:polygain}. The expansion is generated by finite addition, multiplication, and powering of coefficient lists. Equality to \eqref{eq:polyT} follows from distributivity and the finite binomial identities. Alternatively, the checker may accept an expanded list accompanied by exact coefficient equality; an unchecked expansion is never trusted.

All these conditions are bounded rational comparisons. Their conjunction is one PR bit \(\operatorname{cert}(c)\), where \(c\) contains all the finite data. The bit is computable even when a candidate certificate fails. Matrices \(B,C\) need not be supplied by a special mathematical oracle: for example, one may try the inverse of the rational Jacobian at \(x_0\), computed by a bounded pivot procedure, and verify its inverse identity. That choice often works, but completeness of this search is not part of the theorem.

\begin{theorem}[Certified polynomial solver]\label{thm:polysolver}
Assume the checker accepts. Set \(h_0=0\) and \(h_{n+1}=T(h_n)\), and let
\[
 c_0=\max_{i<d}\frac{\abs{T_i(0)}}{v_i},\qquad
 A=\max_{i<d}\sum_{j<d}\abs{C_{ij}}v_j.
\]
For all natural \(n,m\), the following are finite rational inequalities:
\begin{align}
 \abs{(h_n)_i}&\le r v_i &&(i<d),\label{eq:polyball}\\
 \norm{h_{n+1}-h_n}_v&\le c_0 2^{-n},\label{eq:polystep}\\
 \norm{h_{n+m}-h_n}_v&\le 2c_0 2^{-n}(1-2^{-m})\le r2^{-n},\label{eq:polytail}\\
 \norm{F(x_0+h_n)}_\infty&\le A c_0 2^{-n}.\label{eq:polyresidual}
\end{align}
The iteration and these bounds are uniform in the numerical polynomial code, dimension, and precision. Their proofs reduce to QF induction over explicitly coded finite arrays.
\end{theorem}
\begin{proof}
By \eqref{eq:polylip} and \eqref{eq:polycert}, \(\norm{T(h)-T(k)}_v\le\tfrac12\norm{h-k}_v\). In particular, if \(\norm h_v\le r\), then
\[
 \norm{T(h)}_v\le\norm{T(0)}_v+\tfrac12\norm h_v
 \le r/2+r/2=r.
\]
This proves the invariant ball by induction. The first successive difference has norm \(c_0\). Each subsequent difference is at most half the preceding one, proving \eqref{eq:polystep}. Sum the \(m\) consecutive differences beginning at \(n\). The finite geometric identity gives \eqref{eq:polytail}; after multiplying by \(2^{n+m}\), this is an ordinary rational inequality with no infinite summation.

From the defining equation,
\(BF(x_0+h_n)=h_n-h_{n+1}\). Multiplication by the checked left inverse \(C\) yields
\(F(x_0+h_n)=C(h_n-h_{n+1})\). Rowwise triangle inequalities and \eqref{eq:polystep} give \eqref{eq:polyresidual}. Positivity of every denominator is an invariant of the rational representation. Thus every induction step is a finite algebraic implication on the stored state.
\end{proof}

This proves a strong represented consequence. Put \(V=\max(1,\max_{i<d}v_i)\). Choosing
\[
 N(p)=p+\ell\bigl(\max(1,rV,A c_0)\bigr)+2
\]
gives a PR name \(x_0+h_{N(p)}\) whose later values are within \(2^{-p-2}\) in coordinate maximum norm and whose polynomial residual is at most \(2^{-p-2}\). Polynomial evaluation on a bounded box is itself continuous with the explicit bound already proved. Therefore this name is a represented solution of \(F(x)=0\). The object theorem is the universally quantified finite statement \eqref{eq:polyball}--\eqref{eq:polyresidual}, not a quantifier over a new sort of completed real numbers.

\subsection{An automatic radius for Jacobian preconditioning}

The certificate can often be produced directly, without searching over
weights or radii. Let $J$ be the matrix of the linear coefficients of the
expanded polynomial $F(x_0+h)$. Compute and verify $B=J^{-1}$, taking
$C=J$. Then the linear part of $T(h)=h-BF(x_0+h)$ is exactly zero.
After combining equal monomials, write
\[
 T_i(h)=b_i+\sum_{|\alpha|\ge2}t_{i,\alpha}h^\alpha,
 \qquad
 M=\max\left(1,\max_i\sum_{|\alpha|\ge2}
                   |t_{i,\alpha}|\,|\alpha|\right).
\]
Set $v_i=1$ and $r=1/(2M)$. Since $r\le1$, the computed gain obeys
\[
 \sum_j W_{ij}
 =\sum_{|\alpha|\ge2}|t_{i,\alpha}|\,|\alpha|r^{|\alpha|-1}
 \le rM=1/2.
\]
Thus the single remaining numerical test
\[
                 \max_i|b_i|\le 1/(4M)
\]
produces an accepted certificate. All constants are coefficient sums and
one verified rational matrix inverse. This is a quantitative local implicit
equation theorem expressed entirely by finite arithmetic.

For example, with $F(x)=x^2-2$ and $x_0=3/2$, one has $J=3$, $B=1/3$ and
$T(h)=-1/12-h^2/3$. Here $M=1$ and $r=1/2$ pass the test. The explicit
rational iteration therefore constructs the unique positive square root
within $[1,2]$, with the finite tail and residual bounds already proved.
This example requires no bisection oracle or decision of the sign of a
represented real. In a general input, a singular rational $J$ or a failed
residual inequality simply leaves this sufficient method inapplicable.

\subsection{Uniqueness and an a posteriori estimate}

\begin{proposition}[Residual controls distance]\label{prop:polyapost}
For every rational \(z\) with \(\norm z_v\le r\), put \(\eta=\norm{z-T(z)}_v\). Then
\[
 \norm{z-h_n}_v\le r2^{-n}+2\eta\qquad(n\in\N).
\]
In the represented interpretation, there is at most one zero of \(F(x_0+z)\) in the box, and an approximate zero with fixed-point residual \(\eta\) is within \(2\eta\) of the produced name.
\end{proposition}
\begin{proof}
The triangle inequality gives
\(\norm{z-h_{n+1}}_v\le\eta+\tfrac12\norm{z-h_n}_v\).
Starting with \(\norm z_v\le r\), finite induction gives the sharper bound
\(r2^{-n}+2\eta(1-2^{-n})\). Dropping its final nonnegative correction proves the display. A represented exact zero has fixed-point residual as small as any requested rational tolerance; polynomial difference bounds transfer this fact to its rational approximants. Applying the finite estimate at arbitrarily large requested precision proves equality of the names. This supplies the full meaning of uniqueness without a decision procedure for equality of arbitrary real names.
\end{proof}

For the represented uniqueness statement, use rational approximants inside the closed box. Coordinatewise clamping supplies such approximants without increasing coordinate error. Polynomial Lipschitz bounds show that clamping a name of a point in the box preserves its limiting residual. Thus the domain hypothesis used in every finite comparison is explicit.

\subsection{Dependence on parameters and independence of a name}

Let \(T_a\) be a rational polynomial family, with a single invariant box and gain certificate valid for all rational parameters \(a\) in a specified rational parameter box. Suppose the same monomial calculation proves
\[
 \norm{T_a(h)-T_b(h)}_v\le L\norm{a-b}_w
\]
on that box. Then the exact iterations, both starting at zero, satisfy
\begin{equation}\label{eq:paramstability}
 \norm{h_n(a)-h_n(b)}_v\le 2L\norm{a-b}_w.
\end{equation}
Indeed the difference satisfies \(d_{n+1}\le d_n/2+L\norm{a-b}_w\), and its initial value is zero. This is a second finite geometric induction.

If a parameter is itself given by a fixed PR Cauchy name \(a(s)\) in that box, with \(\norm{a(s)-a(t)}_w\le2^{-s}\) for \(s\le t\), define
\[
 b(p)=h_{p+\ell(4r)}\bigl(a(p+\ell(8L))\bigr).
\]
For \(p\le q\), compare \(b(p)\) with the later iteration at the same first parameter, then change the parameter. Equations \eqref{eq:polytail} and \eqref{eq:paramstability} give
\[
 \norm{b(p)-b(q)}_v
 \le r2^{-p-\ell(4r)}+2L2^{-p-\ell(8L)}
 \le 2^{-p-1}.
\]
The shifts also work when \(L=0\). The same proof compares two equivalent input names and proves independence of their particular rational representatives. If the original input approximants can leave the parameter box, coordinatewise rational clamping is a PR nonexpansive map and supplies a convenient invariant-preserving representation.

This is the sense in which the framework is more general than a single approximation algorithm: it supplies composition, perturbation, and extensionality for the actual produced observations, including simultaneous implicit equations. It does not decide whether every polynomial system admits an accepting certificate, and it does not handle a singular Jacobian merely by renaming the problem.

\subsection{Inexact local evaluation without circular error assumptions}

In applications one may replace \(T(h_n)\) by a computed \(\widetilde T_n(h_n)\). The required local error bound must be proved independently of the desired output. Suppose the stronger invariant margin \(\norm{T(0)}_v\le r/4\) holds and every evaluation error is at most \(\delta\le r/4\). Then
\[
 \norm{\widetilde T_n(h)}_v\le r/4+r/2+r/4=r
\]
on the ball. Comparing this approximate iteration with the exact one gives error at most \(2\delta\) after any number of steps. To attain total tolerance \(\epsilon\), choose
\(r2^{-N}\le\epsilon/2\) and
\(\delta\le\min(r/4,\epsilon/4)\).
This derives a local precision demand from a checked stability margin. It does not use the final accuracy statement as its own local premise.

\begin{example}[A genuinely coupled polynomial map]
Consider
\[
 T_1(x,y)=\tfrac1{10}+y+\tfrac{x^2}{100},\qquad
 T_2(x,y)=\tfrac1{20}+\tfrac{x}{32}+\tfrac{y^2}{100}.
\]
Take \(v=(4,1)\) and \(r=1/4\), so the box is \(\abs x\le1\), \(\abs y\le1/4\). Its gain matrix is
\[
 W=\begin{pmatrix}1/50&1\\1/32&1/200\end{pmatrix},\qquad
 Wv=\begin{pmatrix}27/25\\13/100\end{pmatrix}
 \le\tfrac12\begin{pmatrix}4\\1\end{pmatrix}.
\]
The residual at zero is within \((rv)/2=(1/2,1/8)\). Therefore the theorem constructs the unique solution of \((x,y)=T(x,y)\) in the box. The first unweighted row sum is greater than one, so the ordinary unweighted contraction test would fail. Both variables are updated simultaneously; the weighted certificate justifies their mutual dependence.
\end{example}
